% ===================================================================
%  Guia de instalacion - MagikBurger (actualizacion automatica)
%  Compilar en Overleaf con: pdfLaTeX  (archivo en UTF-8)
% ===================================================================
\documentclass[11pt,a4paper]{article}

\usepackage[utf8]{inputenc}
\usepackage[T1]{fontenc}
\usepackage[spanish,es-noshorthands]{babel}
\usepackage{lmodern}
\usepackage{textcomp}
\usepackage[margin=2.3cm]{geometry}
\usepackage{xcolor}
\usepackage{listings}
\usepackage{enumitem}
\usepackage[most]{tcolorbox}
\usepackage{parskip}
\usepackage[hidelinks]{hyperref}

% ---- Estilo de los bloques de comandos ----
\definecolor{codebg}{gray}{0.95}
\lstset{
  basicstyle=\ttfamily\small,
  backgroundcolor=\color{codebg},
  frame=single,
  rulecolor=\color{gray!50},
  framesep=6pt,
  xleftmargin=6pt,
  xrightmargin=6pt,
  breaklines=true,
  breakatwhitespace=false,
  columns=fullflexible,
  keepspaces=true,
  showstringspaces=false,
  aboveskip=8pt,
  belowskip=8pt,
}

% ---- Cajas de aviso ----
\newtcolorbox{importante}{
  colback=red!5!white, colframe=red!70!black,
  title=\textbf{IMPORTANTE}, fonttitle=\bfseries,
  boxrule=0.8pt, arc=2pt, left=6pt, right=6pt, top=4pt, bottom=4pt
}
\newtcolorbox{nota}{
  colback=blue!4!white, colframe=blue!55!black,
  title=\textbf{Nota}, fonttitle=\bfseries,
  boxrule=0.8pt, arc=2pt, left=6pt, right=6pt, top=4pt, bottom=4pt
}

% ---- Atajos para rutas/comandos en linea ----
\newcommand{\rut}[1]{\texttt{#1}}
\newcommand{\bs}{\textbackslash{}}

\setlist[enumerate]{itemsep=4pt, topsep=4pt}

\title{\textbf{Guía paso a paso}\\[4pt]
\large Instalar la versión nueva en el local (UNA sola vez)}
\author{}
\date{}

\begin{document}
\maketitle
\vspace{-1.2cm}

Esta guía es para dejar instalada en la computadora del local la versión con
\textbf{actualización automática}. Después de hacer esto \textbf{una sola vez}, no hay que
volver a instalar nunca más: los cambios futuros llegan solos al abrir la app.

\begin{nota}
Hay dos partes:
\begin{itemize}[leftmargin=1.4em, itemsep=2pt, topsep=2pt]
  \item \textbf{PARTE A:} en TU computadora (la de programar). Genera el instalador.
  \item \textbf{PARTE B:} en la computadora DEL LOCAL. Instala y rescata la base de datos.
\end{itemize}
Hacelas en orden. Tomate tu tiempo: está todo explicado con lo que hay que escribir y dónde tocar.
\end{nota}

% ===================================================================
\section*{PARTE A --- En tu computadora (generar el instalador)}

\subsection*{A.1) Subir el código}
El repositorio donde vive la app necesita tener el código nuevo. En general es:

\begin{lstlisting}
git add -A
git commit -m "Features nuevas + auto-actualizacion"
git push
\end{lstlisting}

\subsection*{A.2) Generar el instalador}
\begin{enumerate}
  \item Abrí la app \textbf{Símbolo del sistema}: tocá el botón \textbf{Inicio} de Windows,
        escribí \rut{cmd} y abrí ``Símbolo del sistema''.
  \item Escribí esto y apretá Enter (cambiá la ruta si tu proyecto está en otro lado):
\begin{lstlisting}
cd "C:\Users\lbera\OneDrive\Desktop\Proyectos\Proyecto Magik\launcher"
\end{lstlisting}
  \item La primera vez, instalá las herramientas (esperá a que termine):
\begin{lstlisting}
npm install
\end{lstlisting}
  \item Generá el instalador:
\begin{lstlisting}
npm run dist
\end{lstlisting}
  \item Cuando termine, el instalador queda en la carpeta
        \rut{...\bs Proyecto Magik\bs launcher\bs dist\bs}.
        Buscá un archivo que se llame algo como
        \textbf{\rut{MagikBurger-Launcher-Setup-1.0.1.exe}}.
  \item Copiá ese \rut{.exe} a un pendrive (o mandalo a la PC del local como prefieras).
\end{enumerate}

% ===================================================================
\section*{PARTE B --- En la computadora del local (instalar SIN perder datos)}

\begin{importante}
Primero rescatamos la base de datos actual (donde están todos los pedidos y la
configuración) para que no se pierda nada.
\end{importante}

\subsection*{B.1) Cerrar la app actual}
\begin{enumerate}
  \item Si la app de MagikBurger está abierta, \textbf{cerrala} (la X de arriba a la derecha).
\end{enumerate}

\subsection*{B.2) Encontrar y copiar la base de datos actual (el rescate)}
\begin{enumerate}
  \item Apretá las teclas \textbf{Windows + R} (se abre una ventanita ``Ejecutar'').
  \item Escribí exactamente esto y apretá Enter:
\begin{lstlisting}
%LOCALAPPDATA%\Programs
\end{lstlisting}
        (Se abre el explorador de archivos. Ahí suele estar instalada la app.)
  \item Buscá una carpeta llamada \textbf{\rut{MagikBurger Launcher}} (o parecida) y entrá.
  \item Entrá a: \textbf{\rut{resources}} $\rightarrow$ \textbf{\rut{backend\_app}}.
  \item Ahí adentro vas a ver un archivo llamado \textbf{\rut{magikburger.db}}.
  \item \textbf{Copialo} (clic derecho $\rightarrow$ Copiar) y \textbf{pegalo en el Escritorio}
        (clic derecho en el Escritorio $\rightarrow$ Pegar). Ese es el rescate: NO lo borres.
\end{enumerate}

\begin{nota}
Si en \rut{\%LOCALAPPDATA\%\bs Programs} no aparece la carpeta, probá también en
\rut{C:\bs Program Files\bs MagikBurger Launcher\bs resources\bs backend\_app\bs}.
El archivo a copiar siempre es \textbf{\rut{magikburger.db}}.
\end{nota}

\subsection*{B.3) Instalar la versión nueva}
\begin{enumerate}
  \item Hacé doble clic en el instalador (\rut{MagikBurger-Launcher-Setup-...exe}) que trajiste.
  \item Seguí los pasos (Siguiente / Instalar). Si Windows muestra un aviso azul de
        ``Windows protegió tu PC'', tocá \textbf{Más información} $\rightarrow$
        \textbf{Ejecutar de todas formas}.
  \item Cuando termine, \textbf{abrí la app una vez} (que cargue) y después \textbf{cerrala}.
        (Esto crea la carpeta donde va a vivir la base de ahora en más.)
\end{enumerate}

\subsection*{B.4) Poner la base rescatada en su lugar definitivo}
\begin{enumerate}
  \item Apretá de nuevo \textbf{Windows + R}.
  \item Escribí exactamente esto y apretá Enter:
\begin{lstlisting}
%APPDATA%\MagikBurger Launcher\data
\end{lstlisting}
        (Se abre una carpeta llamada \textbf{\rut{data}}. Adentro va a haber un
        \rut{magikburger.db} ``vacío'', el que creó la app recién.)
  \item Copiá el \textbf{\rut{magikburger.db} del Escritorio} (el rescatado en el paso B.2) y
        \textbf{pegalo dentro de esta carpeta \rut{data}}, eligiendo
        \textbf{``Reemplazar el archivo de destino''} cuando pregunte.
\end{enumerate}

\subsection*{B.5) Abrir y verificar}
\begin{enumerate}
  \item Abrí la app de nuevo.
  \item Fijate que estén \textbf{los pedidos y datos de siempre} (productos, repartidores, etc.).
        Si está todo: ¡listo! Quedó instalada y con sus datos reales.
\end{enumerate}

% ===================================================================
\section*{¡Y eso es todo!}

A partir de ahora, cada cambio que hagas y subas se actualiza solo en el local
\textbf{la próxima vez que abran la app}. No hay que volver a instalar nada ni tocar la base.

\subsection*{Si algo sale mal}
\begin{itemize}[leftmargin=1.4em]
  \item La base rescatada del Escritorio sigue ahí: podés repetir el paso B.4.
  \item En \rut{\%APPDATA\%\bs MagikBurger Launcher\bs} hay un archivo \rut{launcher.log}
        que registra qué pasó en cada arranque (sirve para diagnosticar).
\end{itemize}

\end{document}
